\chapter{实验平台与成本优化方法}
\label{chap:methods}

本章介绍本文实验的整体平台设计与七种成本优化方法的具体实现。
第~\ref{sec:platform}~节描述实验环境、数据集与评估流程；
第~\ref{sec:gateway}~节介绍本地 token 计量网关的设计；
第~\ref{sec:baseline}~至第~\ref{sec:hybrid}~节依次描述各方法的实现原理
与关键参数；最后第~\ref{sec:method-summary}~节对所有方法进行统一汇总。

%% ─────────────────────────────────────────────────────────────────────────────
\section{实验平台设计}
\label{sec:platform}

\subsection{软硬件环境}
\label{subsec:hw-env}

本文所有实验在单台 Azure 云服务器上运行，主要软硬件配置如表~\ref{tab:hw-env}~所示。

\begin{table}[htbp]
  \centering
  \caption{实验硬件与软件环境}
  \label{tab:hw-env}
  \begin{tabular}{ll}
    \toprule
    项目 & 配置 \\
    \midrule
    操作系统 & Ubuntu 22.04 LTS \\
    CPU & 16 vCPU（Intel Xeon） \\
    内存 & 64 GB \\
    容器运行时 & Docker 24.0，官方 SWE-bench 镜像 \\
    Python & 3.11（Anaconda） \\
    LLM 模型 & MiniMaxAI/MiniMax-M1（通过 SciTix API） \\
    网关 & 本地 FastAPI 透明代理，端口 18110 \\
    \bottomrule
  \end{tabular}
\end{table}

\subsection{数据集}
\label{subsec:dataset}

本文使用 \textbf{SWE-bench Verified} 子集，共 500 条经人工验证的高质量 instance，
覆盖 \texttt{astropy}、\texttt{django}、\texttt{sympy}、\texttt{matplotlib}、
\texttt{scikit-learn} 等主流 Python 开源仓库。
每条 instance 由仓库快照、issue 描述和参考补丁三部分构成（见第~\ref{subsec:swebench-instance}~节）。
所有 7 种方法均在同一组 500 条 instance 上独立运行，保证结果的可比性。

\subsection{实验执行流程}
\label{subsec:run-pipeline}

整体执行流程见算法~\ref{alg:run-pipeline}。

\begin{algorithm}[H]
\caption{批量实验执行流程}
\label{alg:run-pipeline}
\KwIn{实例集合 $\mathcal{D}$（500 条），方法列表 $\mathcal{M}$，最大步数 $K=40$}
\KwOut{每方法的 resolve rate、avg tokens、CPR}

\BlankLine
\For{每种方法 $m \in \mathcal{M}$}{
  \For{每个实例 $\mathcal{T}_i = (r_i, s_i) \in \mathcal{D}$}{
    $\hat{s}_i \leftarrow \text{BuildContext}(s_i,\; r_i,\; m)$ \tcp*{按方法构建 task prompt}
    拉取 Docker 镜像 $\text{img}(r_i)$，挂载 \texttt{/testbed}\;
    将 \texttt{OPENAI\_BASE\_URL} 指向本地网关（端口 18110）\;
    $\hat{p}_i \leftarrow \text{MiniSWEAgent}(\hat{s}_i,\; K)$ \tcp*{运行 agent，至多 K 步}
    $\text{resolved}_i \leftarrow \text{RunTests}(\hat{p}_i,\; \mathcal{T}_i)$ \tcp*{在 Docker 中执行测试}
    保存 \texttt{trajectory.json}，\texttt{result.json}，\texttt{patch.diff}\;
  }
  聚合统计：$\text{ResolveRate}(m)$，$\overline{C}(m)$，$\overline{K}(m)$，$\text{CPR}(m)$\;
}
\Return 所有方法的评估指标矩阵\;
\end{algorithm}

每个 instance 独立在 Docker 容器中运行，互不干扰。
所有方法共享相同的 mini-swe-agent 执行器和相同的测试评估逻辑，
区别仅在于 \texttt{BuildContext} 阶段注入的 task prompt 不同。

\subsection{评估指标}
\label{subsec:metrics}

遵循第~\ref{sec:cost-capability-tradeoff}~节的成本--能力框架，使用以下四个指标：

\begin{itemize}
  \item \textbf{Resolve Rate}：正确解决的 instance 比例（FAIL\_TO\_PASS 测试全部通过）
  \item \textbf{Avg Tokens}：所有 instance 的平均 \texttt{total\_tokens}（input + output）
  \item \textbf{Avg Steps}：平均执行步数
  \item \textbf{CPR}（Cost Per Resolved）：$\overline{C}_{\text{total}} \;/\; \text{ResolveRate}$，
        综合反映成本与能力的联合效率
\end{itemize}


%% ─────────────────────────────────────────────────────────────────────────────
\section{本地 Token 计量网关}
\label{sec:gateway}

\subsection{网关设计动机}
\label{subsec:gateway-motivation}

mini-swe-agent 通过 OpenAI-compatible API 调用 LLM。
为精确计量每一步的 token 消耗，而不依赖上游服务商的计费日志，
本文在 agent 与模型 API 之间插入一个\textbf{本地透明代理网关}（图~\ref{fig:gateway-arch}）。

\begin{algorithm}[H]
\caption{【占位图】此处应为本地 Token 计量网关架构图}
\label{fig:gateway-arch}
\tcc{图示内容：mini-swe-agent → 本地网关（FastAPI, :18110）→ 上游 SciTix API}
\tcc{网关拦截请求/响应，从 usage 字段提取 token 数并写入 JSONL 日志}
\tcc{同时将实验名、方法名、instance\_id 附加到每条日志记录}
\end{algorithm}

\subsection{网关实现}
\label{subsec:gateway-impl}

网关基于 \texttt{FastAPI} 实现，核心逻辑如下：

\begin{lstlisting}[style=promptstyle, caption={网关核心：拦截 token 用量并落盘（节选自 src/gateway/server.py）}, label={lst:gateway-core}]
@app.post("/v1/chat/completions")
async def chat_completions(request: Request) -> JSONResponse:
    body = await request.json()
    instance_id = request.headers.get("X-Instance-Id", "")
    strategy    = request.headers.get("X-Strategy-Name", "")

    # 转发至上游 API
    async with httpx.AsyncClient(timeout=180.0) as client:
        resp = await client.post(
            f"{upstream_url}/chat/completions",
            json=body,
            headers={"Authorization": f"Bearer {upstream_key}"},
        )

    data  = resp.json()
    usage = data.get("usage", {})

    # 记录 token 用量到 JSONL
    record = {
        "ts":           time.time(),
        "instance_id":  instance_id,
        "strategy":     strategy,
        "input_tokens":  usage.get("prompt_tokens",     0),
        "output_tokens": usage.get("completion_tokens", 0),
        "total_tokens":  usage.get("prompt_tokens", 0)
                         + usage.get("completion_tokens", 0),
    }
    with token_log_file.open("a") as f:
        f.write(json.dumps(record) + "\n")

    return JSONResponse(content=data)
\end{lstlisting}

启动时通过 YAML 配置文件指定上游 API 地址和模型名，
agent 侧只需将 \texttt{OPENAI\_BASE\_URL=http://127.0.0.1:18110/v1} 即可无感接入：

\begin{lstlisting}[style=promptstyle, caption={网关配置文件示例（configs/gateway.scitix-minimax.yaml）}, label={lst:gateway-config}]
upstream_url: "${SCITIX_API_BASE_URL}"
upstream_key: "${SCITIX_API_KEY}"
default_model: "MiniMaxAI/MiniMax-M1"
log_dir: "logs/gateway_scitix_minimax"
host: "127.0.0.1"
port: 18110
timeout: 180
\end{lstlisting}

每一步的 token 记录最终汇总写入 \texttt{trajectory.json} 的 \texttt{total\_usage} 字段，
作为后续分析的唯一 token 数据来源。


%% ─────────────────────────────────────────────────────────────────────────────
\section{方法一：baseline\_raw（无压缩基线）}
\label{sec:baseline}

\subsection{方法描述}

\texttt{baseline\_raw} 是无任何增强的对照基线。
它直接将原始 \texttt{problem\_statement} 作为 task prompt 传给 mini-swe-agent，
不增加任何检索结果、工作流提示或 skill memory。

\begin{algorithm}[H]
\caption{baseline\_raw：无压缩基线}
\label{alg:baseline}
\KwIn{问题描述 $s$，仓库快照 $r$}
\KwOut{task prompt $\hat{s}$}
\BlankLine
$\hat{s} \leftarrow s$ \tcp*{直接使用原始 issue，不做任何增强}
\Return $\hat{s}$\;
\end{algorithm}

\subsection{参数设置}

\begin{table}[htbp]
  \centering
  \caption{baseline\_raw 参数配置}
  \label{tab:baseline-params}
  \begin{tabular}{lll}
    \toprule
    参数 & 值 & 说明 \\
    \midrule
    \texttt{method.name} & \texttt{none} & 不启用任何上下文增强 \\
    \texttt{agent.max\_steps} & 40 & 最大执行步数 \\
    \bottomrule
  \end{tabular}
\end{table}

\subsection{预期行为}

agent 在空白上下文中自行决定搜索策略，通常需要多轮 \texttt{grep}/\texttt{find}
才能定位目标文件，搜索阶段产生大量 token。
该方法的价值不在于节省成本，而在于提供对照组，用于衡量其他方法的增益。

%% ─────────────────────────────────────────────────────────────────────────────
\section{方法二：rag\_topk（文件级 RAG 检索）}
\label{sec:rag-topk}

\subsection{方法描述}

\texttt{rag\_topk} 在 agent 运行前，先根据 issue 文本从仓库中检索出少量候选文件，
并将这些文件的代码片段注入 task prompt。
其核心是：\textbf{用轻量规则检索替代 agent 的早期搜索阶段}，
减少 agent 在全仓库范围执行 \texttt{grep}/\texttt{find} 的 token 开销。

\subsection{实现流程}

完整流程见算法~\ref{alg:rag-topk}。

\begin{algorithm}[H]
\caption{rag\_topk：文件级 BM25 检索注入}
\label{alg:rag-topk}
\KwIn{问题描述 $s$，仓库快照 $r$（Docker \texttt{/testbed}），$k=4$，\texttt{snippet\_lines}=220}
\KwOut{增强后的 task prompt $\hat{s}$}

\BlankLine
\tcc{Step 1: 从 issue 中提取查询词}
$T \leftarrow \text{ExtractTerms}(s)$ \tcp*{提取标识符、类名、函数名等，去除停用词，取前 12 个}

\BlankLine
\tcc{Step 2: 在 Docker 中用 rg 检索候选文件}
$\text{cmd} \leftarrow$ \texttt{"rg -l -i -m 1 <pattern> . --glob '!*.json' | head -24"}\;
$\text{candidates} \leftarrow \text{DockerRun}(r,\; \text{cmd})$ \tcp*{最多 24 个候选文件路径}

\BlankLine
\tcc{Step 3: 读取候选文件前 snippet\_lines 行}
\For{每个候选文件 $f$}{
  $\text{content}[f] \leftarrow \text{DockerRun}(r,\; \texttt{"sed -n '1,220p' "} + f)$\;
}

\BlankLine
\tcc{Step 4: BM25 相关性打分，选 top-k 文件}
\For{每个候选文件 $f$}{
  $\text{score}[f] \leftarrow \text{BM25}(T,\; \text{content}[f]) + 2 \times \mathbb{1}[T \cap \text{path}(f) \neq \varnothing]$\;
}
$F_k \leftarrow \text{TopK}(\text{score},\; k=4)$\;

\BlankLine
\tcc{Step 5: 拼接候选文件片段到 prompt}
$\hat{s} \leftarrow s + \text{"Repository hints:"} + \bigcup_{f \in F_k} \text{"[Candidate File] "} + f + \text{content}[f]$\;
$\hat{s} \leftarrow \hat{s} + \text{"The additional context is a starting point..."}$\;
\Return $\hat{s}$\;
\end{algorithm}

\subsection{BM25 打分细节}

本文使用轻量自实现的 BM25（不依赖外部索引库），
查询词集合来自对 issue 文本的正则提取，并对文件路径中的关键词匹配给予额外加权：

\begin{lstlisting}[style=promptstyle, caption={TopKFileSelector BM25 打分（节选自 src/compression/context\_compressor.py）}, label={lst:bm25}]
def _bm25_score(query, files):
    query_terms = re.findall(r"\w+", query.lower())
    scores = {}
    for path, content in files.items():
        words = re.findall(r"\w+", content.lower())
        word_freq = {w: words.count(w) for w in set(words)}
        score = sum(word_freq.get(t, 0) for t in query_terms)
        # 路径匹配额外加权
        path_score = sum(2 for t in query_terms if t in path.lower())
        scores[path] = score + path_score
    return scores
\end{lstlisting}

\subsection{参数设置}

\begin{table}[htbp]
  \centering
  \caption{rag\_topk 参数配置}
  \label{tab:ragtopk-params}
  \begin{tabular}{lll}
    \toprule
    参数 & 值 & 说明 \\
    \midrule
    \texttt{top\_k\_files} & 4 & 注入候选文件数量 \\
    \texttt{snippet\_lines} & 220 & 每个候选文件的最大行数截取 \\
    \texttt{max\_candidate\_files} & 24 & \texttt{rg} 检索的最大候选数 \\
    \texttt{agent.max\_steps} & 40 & 最大执行步数 \\
    \bottomrule
  \end{tabular}
\end{table}


%% ─────────────────────────────────────────────────────────────────────────────
\section{方法三：rag\_function（函数级裁剪检索）}
\label{sec:rag-function}

\subsection{方法描述}

\texttt{rag\_function} 在 \texttt{rag\_topk} 的基础上增加了\textbf{函数级裁剪}步骤：
对 top-k 候选文件，按 Python 顶层 \texttt{def}/\texttt{class} 边界切分代码块，
只保留与 issue 相关的函数和类，丢弃无关代码块，从而进一步减少注入的 token 数。

\subsection{函数级裁剪实现}

\begin{algorithm}[H]
\caption{FunctionLevelTrimmer：函数级代码裁剪}
\label{alg:func-trim}
\KwIn{文件内容 $C$，查询词集合 $T$}
\KwOut{裁剪后的文件内容 $C'$}

\BlankLine
$\text{blocks} \leftarrow \text{SplitByTopLevelDef}(C)$ \tcp*{按顶层 def/class 切分为块}
$T_{\text{exp}} \leftarrow \text{ExpandIdentifiers}(T)$ \tcp*{展开驼峰和下划线标识符}
$C' \leftarrow []$\;

\BlankLine
\For{每个代码块 $(t, \text{lines}) \in \text{blocks}$}{
  \uIf{$t = \texttt{import}$}{
    保留（import 语句始终保留）\;
  }
  \uElseIf{$t = \texttt{class\_sig}$}{
    保留类签名行\;
  }
  \uElseIf{$T_{\text{exp}} \cap \text{Terms}(\text{lines}) \neq \varnothing$}{
    保留（函数/类与查询词有交集）\;
  }
  \Else{
    丢弃（无关代码块）\;
  }
}
\Return $C'$\;
\end{algorithm}

标识符展开（\texttt{ExpandIdentifiers}）将驼峰命名和下划线命名拆分成子词，
例如 \texttt{separability\_matrix} 展开为 \{\texttt{separability}, \texttt{matrix}\}，
提高跨命名风格的匹配召回率。

\subsection{参数设置与风险}

\begin{table}[htbp]
  \centering
  \caption{rag\_function 参数配置}
  \label{tab:ragfunc-params}
  \begin{tabular}{lll}
    \toprule
    参数 & 值 & 说明 \\
    \midrule
    \texttt{top\_k\_files} & 4 & 同 rag\_topk \\
    \texttt{snippet\_lines} & 220 & 先截取前 220 行，再函数级裁剪 \\
    \texttt{enable\_function\_trim} & True & 启用函数级裁剪 \\
    \texttt{agent.max\_steps} & 40 & 最大执行步数 \\
    \bottomrule
  \end{tabular}
\end{table}

该方法的风险在于：若裁剪过激，可能丢失修复所需的辅助函数、类属性或导入关系，
导致 agent 在后续执行中仍需重新读取文件，抵消压缩效果。

%% ─────────────────────────────────────────────────────────────────────────────
\section{方法四：llmlingua\_original（语义文件重排序）}
\label{sec:llmlingua}

\subsection{方法描述}

\texttt{llmlingua\_original} 将 LLMLingua~\cite{jiang2023llmlingua} 的语义选择思想
适配到文件级检索场景。它不是对文件内容进行 token 级压缩，而是用
LLMLingua 的小型语言模型（\texttt{deepseek-coder-1.3b-base}）对候选文件进行
\textbf{语义相关性重排序}，以替代 BM25 的词频统计。

\subsection{LLMLingua 适配实现}

\begin{algorithm}[H]
\caption{llmlingua\_original：语义文件重排序}
\label{alg:llmlingua}
\KwIn{问题描述 $s$，候选文件集合 $F$（来自 rg 检索），$k=4$}
\KwOut{增强后的 task prompt $\hat{s}$}

\BlankLine
\tcc{Step 1: 用 rg 和关键词召回候选文件（同 rag\_topk）}
$F \leftarrow \text{CollectCandidates}(s, r)$\;

\BlankLine
\tcc{Step 2: 用 LLMLingua 对候选文件重排序}
$\text{compressor} \leftarrow \text{PromptCompressor}(\texttt{"deepseek-coder-1.3b"})$\;
\For{每个文件 $f \in F$}{
  将 $f$ 包装为 \texttt{SimpleNamespace(path=f, code\_content=content)}\;
}
$F_k \leftarrow \text{compressor.select\_cross\_files}(F,\; \text{question}=s,\; \text{topk}=4)$\;

\BlankLine
\tcc{Step 3: 同时启用函数级裁剪}
$F_k \leftarrow \{f: \text{FunctionTrim}(\text{content}[f], s) \;\text{for}\; f \in F_k\}$\;

\BlankLine
$\hat{s} \leftarrow s + \text{format}(F_k)$\;
\Return $\hat{s}$\;
\end{algorithm}

LLMLingua 的 \texttt{select\_cross\_files} 方法通过计算文件内容在问题条件下的
困惑度（perplexity）来评估文件的语义相关性，
相比 BM25 更能捕捉语义层面的关联，而非简单关键词共现。

\subsection{参数设置}

\begin{table}[htbp]
  \centering
  \caption{llmlingua\_original 参数配置}
  \label{tab:llmlingua-params}
  \begin{tabular}{lll}
    \toprule
    参数 & 值 & 说明 \\
    \midrule
    \texttt{llmlingua\_model\_name} & \texttt{deepseek-coder-1.3b-base} & 重排序用小模型 \\
    \texttt{llmlingua\_device\_map} & \texttt{cpu} & 在 CPU 上运行（避免 GPU 竞争） \\
    \texttt{top\_k\_files} & 4 & 保留文件数 \\
    \texttt{enable\_function\_trim} & True & 启用函数级裁剪 \\
    \texttt{agent.max\_steps} & 40 & 最大执行步数 \\
    \bottomrule
  \end{tabular}
\end{table}


%% ─────────────────────────────────────────────────────────────────────────────
\section{方法五：skill\_abstraction（工作流脚手架注入）}
\label{sec:skill-abstraction}

\subsection{方法描述}

\texttt{skill\_abstraction} 属于工作流压缩路径。
它不改变检索逻辑，而是在 task prompt 中额外注入一段\textbf{结构化工作流提示}
（workflow scaffold），引导 agent 按照固定的六步流程解题，
从而减少模型在规划阶段的空转 token。

\subsection{工作流提示内容}

注入的工作流提示为固定字符串，内容如下：

\begin{lstlisting}[style=promptstyle, caption={skill\_abstraction 注入的工作流提示（节选自 src/agent/task\_context.py）}, label={lst:skill-scaffold}]
Structured workflow:
1. Reproduce or understand the issue briefly.
2. Narrow to 1-3 candidate files.
3. Locate the exact function or block to edit.
4. Make the smallest patch that fixes the issue.
5. Verify with a focused check.
6. Submit the final patch only after verification.
\end{lstlisting}

\begin{algorithm}[H]
\caption{skill\_abstraction：工作流脚手架注入}
\label{alg:skill-abstraction}
\KwIn{问题描述 $s$，仓库快照 $r$}
\KwOut{增强后的 task prompt $\hat{s}$}

\BlankLine
\tcc{Step 1: 同 rag\_topk 检索候选文件}
$F_k \leftarrow \text{RetrieveTopK}(s, r, k=4)$\;

\BlankLine
\tcc{Step 2: 构建 prompt 各部分}
$\hat{s} \leftarrow s$\;
$\hat{s} \leftarrow \hat{s} + W_{\text{scaffold}}$ \tcp*{注入结构化工作流提示（固定字符串）}
$\hat{s} \leftarrow \hat{s} + \text{format}(F_k)$ \tcp*{追加候选文件片段}
$\hat{s} \leftarrow \hat{s} + \text{"The additional context is a starting point..."}$\;
\Return $\hat{s}$\;
\end{algorithm}

\subsection{设计意图}

该方法本质上是 trajectory compression：通过抽象一个高层求解流程，
减少 LLM 自行规划时产生的冗余推理步骤。
直觉上，当模型收到明确的步骤指引时，
它可以更快进入"定位文件 → 最小 patch → 验证"的有效路径，
而不是在反复搜索和读取之间徘徊。

\subsection{参数设置}

\begin{table}[htbp]
  \centering
  \caption{skill\_abstraction 参数配置}
  \label{tab:skill-abs-params}
  \begin{tabular}{lll}
    \toprule
    参数 & 值 & 说明 \\
    \midrule
    \texttt{include\_skill\_scaffold} & True & 启用工作流提示注入 \\
    \texttt{top\_k\_files} & 4 & 同时检索 top-k 候选文件 \\
    \texttt{snippet\_lines} & 220 & 文件片段行数 \\
    \texttt{agent.max\_steps} & 40 & 最大执行步数 \\
    \bottomrule
  \end{tabular}
\end{table}

%% ─────────────────────────────────────────────────────────────────────────────
\section{方法六：skill\_memory\_md（静态经验记忆注入）}
\label{sec:skill-memory}

\subsection{方法描述}

\texttt{skill\_memory\_md} 是静态可迁移经验注入方法。
它从根目录的 \texttt{skills.md} 文件中读取跨任务可复用的修复启发式经验，
将其注入 task prompt。与 \texttt{skill\_abstraction} 不同，
这里注入的不是固定流程指令，而是从历史 resolved 样本中人工整理的\textbf{可迁移搜索/编辑/验证启发式规则}。

\subsection{skills.md 内容结构}

\texttt{skills.md} 分为五个小节，分别对应搜索、阅读、编辑、验证和提交五个阶段的启发式规则：

\begin{lstlisting}[style=promptstyle, caption={skills.md 内容节选（搜索与阅读启发式）}, label={lst:skills-md}]
## Search Heuristics
- Search for exact error messages, API names, and test names from the issue.
- If the issue describes argument validation, start near the public API
  entrypoint and trace one level inward.
- If tests are named in the issue metadata, inspect the nearest test module
  to infer target code paths.

## Read Heuristics
- Avoid whole-file dumps unless the file is short.
- Read the function containing the suspected behavior first, then one
  caller and one callee if needed.
- Do not re-read the same file unless you are checking a different symbol.
\end{lstlisting}

\begin{algorithm}[H]
\caption{skill\_memory\_md：静态经验记忆注入}
\label{alg:skill-memory}
\KwIn{问题描述 $s$，仓库快照 $r$，skill 文件路径 \texttt{skills.md}}
\KwOut{增强后的 task prompt $\hat{s}$}

\BlankLine
$F_k \leftarrow \text{RetrieveTopK}(s, r, k=4)$ \tcp*{检索候选文件（同 rag\_topk）}
$M \leftarrow \text{ReadFile}(\texttt{skills.md})$ \tcp*{加载静态 skill memory}

\BlankLine
$\hat{s} \leftarrow s$\;
$\hat{s} \leftarrow \hat{s} + \text{"Transferable skill memory (reusable heuristics):"} + M$\;
$\hat{s} \leftarrow \hat{s} + \text{format}(F_k)$\;
\Return $\hat{s}$\;
\end{algorithm}

\subsection{参数设置}

\begin{table}[htbp]
  \centering
  \caption{skill\_memory\_md 参数配置}
  \label{tab:skill-mem-params}
  \begin{tabular}{lll}
    \toprule
    参数 & 值 & 说明 \\
    \midrule
    \texttt{include\_skill\_memory} & True & 启用 skill memory 注入 \\
    \texttt{skill\_file} & \texttt{skills.md} & skill memory 文件路径 \\
    \texttt{top\_k\_files} & 4 & 同时检索候选文件 \\
    \texttt{snippet\_lines} & 220 & 文件片段行数 \\
    \texttt{agent.max\_steps} & 40 & 最大执行步数 \\
    \bottomrule
  \end{tabular}
\end{table}


%% ─────────────────────────────────────────────────────────────────────────────
\section{方法七：hybrid\_llmlingua（组合压缩）}
\label{sec:hybrid}

\subsection{方法描述}

\texttt{hybrid\_llmlingua} 是本文设计的组合方法，将 LLMLingua 语义重排序、
函数级裁剪和 token budget 截断三个压缩步骤串联，
以最大化上下文压缩力度。它的目标是同时减少注入的文件数量和单文件的 token 长度。

\subsection{三阶段组合压缩流程}

\begin{algorithm}[H]
\caption{hybrid\_llmlingua：三阶段组合压缩}
\label{alg:hybrid}
\KwIn{问题描述 $s$，仓库快照 $r$，$k=4$，\texttt{token\_budget}=2500}
\KwOut{增强后的 task prompt $\hat{s}$}

\BlankLine
\tcc{阶段一：rg 召回候选文件（同 rag\_topk）}
$F \leftarrow \text{CollectCandidates}(s, r)$ \tcp*{最多 24 个候选}

\BlankLine
\tcc{阶段二：LLMLingua 语义重排序，选 top-k}
$F_k \leftarrow \text{LLMLinguaRank}(F,\; s,\; k=4)$\;

\BlankLine
\tcc{阶段三：函数级裁剪}
\For{每个文件 $f \in F_k$}{
  $\text{content}[f] \leftarrow \text{FunctionTrim}(\text{content}[f],\; s)$\;
}

\BlankLine
\tcc{阶段四：token budget 截断}
$F_k \leftarrow \text{TokenBudgetTrim}(F_k,\; \text{budget}=2500)$\;
\tcp*{按文件长度升序依次填入，超出 budget 则截断或跳过}

\BlankLine
$\hat{s} \leftarrow s + \text{format}(F_k)$\;
\Return $\hat{s}$\;
\end{algorithm}

\subsection{token budget 截断逻辑}

\texttt{TokenBudgetTrimmer} 按文件长度（token 数）从短到长依次填入：

\begin{lstlisting}[style=promptstyle, caption={TokenBudgetTrimmer 截断逻辑（节选自 src/compression/context\_compressor.py）}, label={lst:budget-trimmer}]
def trim(self, files, query):
    result, used = {}, 0
    sorted_files = sorted(files.items(),
                          key=lambda kv: count_tokens(kv[1]))
    for path, content in sorted_files:
        tokens    = count_tokens(content)
        remaining = self.token_budget - used
        if tokens <= remaining:
            result[path] = content
            used += tokens
        elif remaining > 200:
            enc       = tiktoken.get_encoding("cl100k_base")
            truncated = enc.decode(enc.encode(content)[:remaining])
            result[path] = truncated + "\n# [TRUNCATED BY TOKEN BUDGET]"
            used += remaining
        if used >= self.token_budget:
            break
    return result
\end{lstlisting}

\subsection{参数设置}

\begin{table}[htbp]
  \centering
  \caption{hybrid\_llmlingua 参数配置}
  \label{tab:hybrid-params}
  \begin{tabular}{lll}
    \toprule
    参数 & 值 & 说明 \\
    \midrule
    \texttt{top\_k\_files} & 4 & LLMLingua 保留文件数 \\
    \texttt{snippet\_lines} & 220 & 初始截取行数 \\
    \texttt{enable\_function\_trim} & True & 启用函数级裁剪 \\
    \texttt{token\_budget} & 2500 & token budget 上限（tiktoken cl100k\_base） \\
    \texttt{llmlingua\_model\_name} & \texttt{deepseek-coder-1.3b-base} & 重排序模型 \\
    \texttt{agent.max\_steps} & 40 & 最大执行步数 \\
    \bottomrule
  \end{tabular}
\end{table}

\subsection{风险分析}

多重压缩步骤叠加后，存在\textbf{过度压缩}的风险：
若 LLMLingua 重排序遗漏了关键文件，或函数级裁剪丢弃了必要的辅助函数，
再经过 token budget 截断，最终注入的上下文可能已不包含修复所需信息。
此时 agent 仍需自行补充搜索，反而增加了总 token 消耗。
因此该方法需要通过完整实验验证，确认组合压缩是否真正优于单一压缩。


%% ─────────────────────────────────────────────────────────────────────────────
\section{方法总览与对比}
\label{sec:method-summary}

\subsection{统一接口与公平比较}

所有七种方法都遵循第~\ref{subsec:method-io}~节定义的统一接口：
输入为 $(s, r, \theta)$，输出为增强后的 task prompt $\hat{s}$。
mini-swe-agent 的执行器、Docker 环境、token 计量网关、
测试评估逻辑均保持一致，方法间的唯一变量是 \texttt{BuildContext} 阶段的行为。

\subsection{方法参数汇总}

\begin{table}[htbp]
  \centering
  \caption{七种方法的关键参数汇总}
  \label{tab:method-params-all}
  \begin{tabularx}{\textwidth}{lXXXXX}
    \toprule
    方法 & 检索文件数 & 函数裁剪 & LLMLingua & token budget & 工作流/记忆注入 \\
    \midrule
    \texttt{baseline\_raw}       & —    & —    & —    & —    & — \\
    \texttt{rag\_topk}           & 4    & —    & —    & —    & — \\
    \texttt{rag\_function}       & 4    & Yes  & —    & —    & — \\
    \texttt{llmlingua\_original} & 4    & Yes  & Yes  & —    & — \\
    \texttt{skill\_abstraction}  & 4    & —    & —    & —    & Workflow \\
    \texttt{skill\_memory\_md}   & 4    & —    & —    & —    & Memory \\
    \texttt{hybrid\_llmlingua}   & 4    & Yes  & Yes  & 2500 & — \\
    \bottomrule
  \end{tabularx}
\end{table}

\subsection{方法分类与预期优化机制}

\begin{table}[htbp]
  \centering
  \caption{七种方法的优化路径与预期作用阶段}
  \label{tab:method-classification}
  \begin{tabularx}{\textwidth}{llXX}
    \toprule
    方法 & 优化路径 & 主要作用阶段 & 预期节省来源 \\
    \midrule
    \texttt{baseline\_raw}       & 无 & — & 对照组 \\
    \texttt{rag\_topk}           & 上下文压缩 & Search + Read & 减少全仓搜索 \\
    \texttt{rag\_function}       & 上下文压缩 & Search + Read & 减少文件内冗余 \\
    \texttt{llmlingua\_original} & 上下文压缩 & Search + Read & 语义精选高价值文件 \\
    \texttt{skill\_abstraction}  & 工作流压缩 & 全阶段 & 减少规划空转 \\
    \texttt{skill\_memory\_md}   & 工作流压缩 & 全阶段 & 复用跨任务经验 \\
    \texttt{hybrid\_llmlingua}   & 两路结合 & 全阶段 & 同时压文件数和长度 \\
    \bottomrule
  \end{tabularx}
\end{table}

\subsection{方法实现依赖关系}

图~\ref{fig:method-deps}~展示了各方法共享的代码模块依赖关系。

\begin{algorithm}[H]
\caption{【占位图】此处应为方法实现模块依赖关系图}
\label{fig:method-deps}
\tcc{图示内容：各方法 → TaskContextBuilder}
\tcc{TaskContextBuilder 调用：TopKFileSelector / FunctionLevelTrimmer / TokenBudgetTrimmer / LLMLinguaRanker}
\tcc{所有方法共享：AgentRunner → mini-swe-agent → Docker → Token Gateway}
\end{algorithm}

%% ─────────────────────────────────────────────────────────────────────────────
\section{本章小结}
\label{sec:chap05-summary}

本章介绍了完整的实验平台与七种成本优化方法的实现细节，主要内容包括：

\begin{enumerate}
  \item \textbf{实验平台}：基于 Azure 云服务器、官方 SWE-bench Docker 镜像和
        mini-swe-agent 构建统一执行环境，500 条 SWE-bench Verified instance
        在七种方法下独立运行，保证可比性。

  \item \textbf{token 计量网关}：设计本地 FastAPI 透明代理，
        精确记录每步 \texttt{input\_tokens} 和 \texttt{output\_tokens}，
        为成本分析提供统一数据来源。

  \item \textbf{上下文压缩方法}（\texttt{rag\_topk}、\texttt{rag\_function}、
        \texttt{llmlingua\_original}）：通过关键词检索 + BM25 打分 / LLMLingua
        语义重排序 + 函数级裁剪，在 agent 运行前注入精简的候选文件上下文，
        目标是减少搜索和阅读阶段的 token 消耗。

  \item \textbf{工作流压缩方法}（\texttt{skill\_abstraction}、\texttt{skill\_memory\_md}）：
        通过注入结构化工作流提示或静态跨任务经验记忆，引导 agent 更快进入
        有效编辑阶段，减少规划空转。

  \item \textbf{组合压缩方法}（\texttt{hybrid\_llmlingua}）：
        串联语义重排序、函数裁剪和 token budget 截断，追求最大压缩力度，
        同时承担过度压缩的风险。
\end{enumerate}

各方法的实际效果将在第~\ref{chap:results}~章通过完整实验数据进行对比分析。
